\begin{lemma}[Third homology under a band-preserving quotient]
\label{B:band-quotient-H3-surjection}
Let $q:M\twoheadrightarrow N$ be a surjective homomorphism of
finite monoids which restricts to a bijection between their
rectangular-band minimal ideals.  After identifying those ideals
by $q$, the induced maps satisfy
\[
 H_2(BM;k)\xrightarrow{\ \cong\ }H_2(BN;k),\qquad
 H_3(BM;k)\twoheadrightarrow H_3(BN;k)
\]
for every commutative coefficient ring $k$, including $\mathbb Z$.
In particular an actor image with nonzero third homology over a
field cannot be the image of a projective-plane monoid of this
kind, even when the action map is not faithful.
\end{lemma}
\begin{proof}
Write the common band as $I\times J$, and put
$V=\widetilde{kI}$ and $W=\widetilde{kJ}$.
Both action modules are unchanged by $q$: multiplication against
a band element is preserved, and $q$ is injective on the band.
The two augmentation sequences used in the band--Tor argument
are therefore natural under $q$, with identity maps on $V,W$.
Their connecting homomorphisms identify the actual homology map
with
\[
 H_{n+2}(Bq;k)=
 \operatorname{Tor}^{kM}_n(W,V)\longrightarrow
 \operatorname{Tor}^{kN}_n(W,V),\qquad n\geq0.
\]
This is also seen directly from the natural full two-sided bar
comparison in the proof of Lemma~\ref{lem:B-band-Tor}.

In degree zero, the tensor-product relations are identical:
the actions of the source elements on $W,V$ are exactly the
actions of their images, and $q$ is surjective.  Thus the Tor
map in degree zero is an isomorphism.

For degree one, use the ordinary two-sided bar over $k$.
Its first two terms are
\[
 C_1(M)=W\otimes_k kM\otimes_k V,
 \qquad C_0(M)=W\otimes_k V,
\]
with differential
\[
 d(w\otimes m\otimes v)=wm\otimes v-w\otimes mv.
\]
The induced map on $C_0$ is the identity and the map on $C_1$
is surjective.  If $z$ is a one-cycle for $N$, choose any lift
$\widetilde z\in C_1(M)$.  Commutation with the differential
gives $d\widetilde z=dz=0$ in the identical degree-zero module.
Every target cycle therefore lifts to a source cycle, and
boundaries map to boundaries.  This proves the Tor-one
surjection.  All modules in this bar construction are free over
$k$ in the requisite outer factors, so it computes the stated
ordinary Tor groups over any $k$.  Naturality of the two
connecting maps gives the asserted statement for $H_3$.
\end{proof}
